\documentclass[12pt]{article}
% Paquetes básicos primero
\usepackage[utf8]{inputenc}
\usepackage[T1]{fontenc}
\usepackage{ragged2e}
\usepackage{textcomp}
\usepackage{gensymb}
\usepackage{siunitx} 
\usepackage[spanish]{babel}
\usepackage{multicol}

% Geometría y diseño
\usepackage{geometry}
\geometry{
	letterpaper,
	left=2.5cm,
	right=2.5cm,
	top=3cm,
	bottom=3cm,
	headheight=60pt,
	headsep=1cm
}

% Paquetes de formato
\usepackage{setspace}
\usepackage{graphicx}
\usepackage{fancyhdr}
\usepackage{xcolor}
\usepackage{enumitem}
\usepackage{lastpage}
\usepackage{background}
\usepackage{listings}
\usepackage{tcolorbox}
\usepackage{fontawesome}
\usepackage{titlesec}
\usepackage{float}

% Paquetes para referencias
\usepackage[round]{natbib}  % Configuración específica para citas
\bibliographystyle{plainnat}  % Estilo de bibliografía

% URLs y enlaces (hyperref debe ir al final)
\usepackage{url}
\usepackage[hidelinks]{hyperref}

% Configuración de colores personalizados
\definecolor{tecnmBlue}{RGB}{0, 43, 89}
\definecolor{tecnmWhite}{RGB}{255, 255, 255}
\definecolor{codeBackground}{RGB}{240,240,240}
\definecolor{comentarios}{RGB}{63,127,95}
\definecolor{identificador}{RGB}{127,0,85}
\definecolor{cadena}{RGB}{42,0,255}
\definecolor{arduinoBlue}{RGB}{0,129,132}

% Configuración de listings para código Arduino
\lstdefinestyle{arduinoStyle}{
	language=C++,
	backgroundcolor=\color{codeBackground},
	basicstyle=\ttfamily\scriptsize,
	breaklines=true,
	captionpos=b,
	commentstyle=\color{comentarios},
	keywordstyle=\color{identificador},
	stringstyle=\color{cadena},
	numbers=left,
	numberstyle=\tiny\color{gray},
	stepnumber=1,
	numbersep=5pt,
	tabsize=3,
	showspaces=false,
	showstringspaces=false,
	frame=single,
	frameround=tttt,
	framesep=2mm,
	rulecolor=\color{arduinoBlue},
	emph={pinMode,digitalWrite,HIGH,LOW,INPUT,OUTPUT,tone,noTone,delay,delayMicroseconds,millis,INPUT_PULLUP},
	emphstyle={\color{blue}\bfseries},
	morekeywords={setup,loop},
	linewidth=\columnwidth,
	xleftmargin=0.5cm,
	xrightmargin=0.5cm,
	float=h,
	floatplacement=H,
	belowskip=-0.2cm,
	aboveskip=0.2cm
}

% Configuración de cajas personalizadas
\newtcolorbox{componentBox}[2][]{
	colback=tecnmBlue!5,
	colframe=tecnmBlue!40!black,
	title=#2,
	width=\columnwidth,
	#1
}

% Configuración del encabezado y pie de página
\pagestyle{fancy}
\fancyhf{}
\renewcommand{\headrulewidth}{0pt}
\renewcommand{\footrulewidth}{0pt}

% Definir comandos para reducir la carga
\newcommand{\smallgray}[1]{{\scriptsize\textcolor{gray}{#1}}}

\fancyhead[L]{%
	\hspace*{0.5cm}\raisebox{-2.2\height}{\includegraphics[height=1cm]{images/logotecnm.png}}%
}
\fancyhead[R]{%
	\raisebox{-1.7\height}{%
		\begin{tabular}[b]{r}
			{\scriptsize\textbf{Instituto Tecnológico de Morelia}}\\[-0.6em]
			\smallgray{Subdirección Académica}\\[-0.6em]
			\smallgray{Departamento de Sistemas y Computación}\\[-0.6em]
			\smallgray{Laboratorio de Sistemas Programables}\\[-0.4em]
			{\scriptsize Página \thepage\ de \pageref{LastPage}}
		\end{tabular}%
	}%
}

\makeatletter
\renewcommand\@biblabel[1]{#1.}
\makeatother

\fancyfoot[L]{%
	\small
	Implementación de un Buzzer
	
	\vspace{0.1cm}
	
	\small\textcolor{gray}{Plantilla de Reporte de Laboratorio}
}
\fancyfoot[R]{%
	\small
	\textbf{Equipo 1}%
}

% Definir comando para citas
\newcommand{\cita}[1]{\cite{#1}}

% Estilo de títulos
\titleformat{\section}
{\color{tecnmBlue}\Large\bfseries}
{\thesection}{1em}{}

\titleformat{\subsection}
{\color{tecnmBlue}\large\bfseries}
{\thesubsection}{1em}{}

% Configuración de siunitx
\sisetup{
	detect-all,
	detect-family,
	per-mode=symbol,
	range-phrase=--,
	range-units=single,
	output-decimal-marker={,}
}

\begin{document}
	\pagenumbering{arabic}
	
	% Portada (una columna)
	% portada.tex
% Configuración del fondo solo para la portada
\backgroundsetup{
	scale=1,
	angle=0,
	opacity=0.2,
	contents={\includegraphics[width=\paperwidth,height=\paperheight]{images/fondo_portada.png}}
}
\BgThispage

\begin{center}
	\vspace*{2cm}
	
	\Large\textbf{Transistor BJT y MOSFET}
	
	\vspace{1.5cm}
	
	Autor(es):
	
	\vspace{0.5cm}
	
	Mario Eduardo Sánchez Mejía\\
	Fidel Alberto Zarco Áviles
	
	\vspace{0.3cm}
	
	\small{21120721@morelia.tecnm.mx}\\
	\small{l20121258@morelia.tecnm.mx}
	
	\vspace{1.5cm}
	
	Asesor(@s):
	
	\vspace{0.5cm}
	
	Luis Ulises Chávez Campos
	
	\vspace{0.5cm}
	
	\begin{center}
		\begin{abstract}
			\noindent
			\justifying
			Se presenta el desarrollo e implementación de un sistema de control electrónico basado en transistores como interruptores, utilizando microcontroladores Arduino. La implementación incorpora componentes específicos como el transistor BD235 y el MOSFET IRLZ14, junto con potenciómetros, para gestionar cargas de alta potencia, particularmente bombillas de 12V. El desarrollo experimental facilita la comprensión práctica del comportamiento de transistores en aplicaciones de conmutación, además de demostrar la implementación efectiva de sistemas de control mediante señales tanto digitales como analógicas. La integración resultante evidencia la aplicación práctica de principios fundamentales en electrónica de potencia y programación de microcontroladores en sistemas embebidos.
			
			\vspace{0.5cm}
			\noindent
			\textbf{Palabras clave:} Arduino, Transistores, MOSFET, Sistemas de Control, Electrónica de Potencia
		\end{abstract}
	\end{center}
\end{center}

% Desactivar el fondo para las siguientes páginas
\clearpage
\backgroundsetup{contents={}}
	% Inicio de dos columnas para el contenido principal
	\begin{multicols}{2}
		\setlength{\columnsep}{1cm} % Espacio entre columnas
		\tableofcontents
		\listoffigures
		\section{Introducción}
La presente documentación técnica establece los fundamentos teóricos y especificaciones técnicas de los componentes semiconductores y elementos pasivos utilizados en sistemas de control de cargas mediante transistores. El análisis abarca desde los principios fundamentales de la teoría de circuitos hasta las características específicas de los dispositivos semiconductores \citep{malvino2021electronic}.

\subsection{Transistor BJT BD235}
\begin{componentBox}{Especificaciones Técnicas y Configuración \citep{onsemi2021bd235}}
	\begin{itemize}[leftmargin=*,itemsep=1pt,parsep=1pt]
		\item \textbf{Clasificación}: Transistor Bipolar de Unión NPN de Potencia
		\item \textbf{Parámetros Característicos}:
		\begin{itemize}[itemsep=0pt,parsep=0pt]
			\item Tensión colector-emisor máxima (Vceo): \SI{50}{\volt}
			\item Corriente de colector máxima (Ic): \SI{2}{\ampere}
			\item Factor de amplificación de corriente (hFE): 40-250
			\item Configuración física: Encapsulado TO-126
		\end{itemize}
		\item \textbf{Configuración de Terminales}:
		\begin{itemize}[itemsep=0pt,parsep=0pt]
			\item Terminal de Colector (C): Electrodo principal de salida
			\item Terminal de Base (B): Electrodo de control
			\item Terminal de Emisor (E): Electrodo de referencia
		\end{itemize}
	\end{itemize}
\end{componentBox}

\subsection{MOSFET IRLZ44}
\begin{componentBox}{Características Técnicas y Operación \citep{infineon2020irlz44}}
	\begin{itemize}[leftmargin=*,itemsep=1pt,parsep=1pt]
		\item \textbf{Parámetros Operativos}:
		\begin{itemize}[itemsep=0pt,parsep=0pt]
			\item Tensión drenador-fuente máxima (Vds): \SI{55}{\volt}
			\item Corriente de drenador máxima (Id): \SI{47}{\ampere}
			\item Resistencia en conducción (Rds(on)): \SI{0.022}{\ohm}
			\item Tensión umbral de compuerta (Vgs(th)): \SI{1}{\volt}-\SI{2}{\volt}
		\end{itemize}
		\item \textbf{Estructura Terminal}:
		\begin{itemize}[itemsep=0pt,parsep=0pt]
			\item Terminal de Compuerta (G): Control de conducción
			\item Terminal de Drenador (D): Conducción principal
			\item Terminal de Fuente (S): Referencia de potencial
		\end{itemize}
	\end{itemize}
\end{componentBox}

\subsection{Potenciómetro de Precisión \SI{10}{\kilo\ohm}}
\begin{componentBox}{Parámetros Técnicos \citep{horowitz2015art}}
	\begin{itemize}[leftmargin=*,itemsep=1pt,parsep=1pt]
		\item \textbf{Especificaciones Fundamentales}:
		\begin{itemize}[itemsep=0pt,parsep=0pt]
			\item Resistencia nominal: \SI{10}{\kilo\ohm}
			\item Configuración mecánica: Sistema rotativo
			\item Característica de variación: Respuesta lineal
		\end{itemize}
		\item \textbf{Disposición de Terminales}:
		\begin{itemize}[itemsep=0pt,parsep=0pt]
			\item Terminal inicial: Punto de referencia
			\item Terminal variable: Cursor de control
			\item Terminal final: Límite resistivo
		\end{itemize}
	\end{itemize}
\end{componentBox}

\subsection{Sistema de Iluminación \SI{12}{\volt}}
\begin{componentBox}{Parámetros Operativos \citep{boylestad2013electronic}}
	\begin{itemize}[leftmargin=*,itemsep=1pt,parsep=1pt]
		\item \textbf{Especificaciones Eléctricas}:
		\begin{itemize}[itemsep=0pt,parsep=0pt]
			\item Tensión nominal de operación: \SI{12}{\volt}
			\item Potencia nominal: \SI{10}{\watt}
			\item Corriente de operación: \SI{0.83}{\ampere}
		\end{itemize}
	\end{itemize}
\end{componentBox}
		%\input{sections/marcoTeorico}
		%\input{sections/metodologia}
	\end{multicols}
	
	% Referencias (una columna)
	\clearpage
	\addcontentsline{toc}{section}{Referencias}
	\bibliography{referencias}
	
\end{document}